\begin{thema}{Β}
Δίνεται ο παρακάτω \textbf{πίνακας προσήμων} της $f'(x)$ για μια συνεχή $f:[-3, 5]\to\mathbb{R}$:

\begin{center}
\begin{tabular}{|c|c|c|c|c|c|}
\hline
$x$ & $-3$ & $\cdots$ & $0$ & $\cdots$ & $5$ \\
\hline
$f'(x)$ & & $+$ & $0$ & $-$ & \\
\hline
$f(x)$ & $-2$ & $\nearrow$ & $4$ & $\searrow$ & $1$ \\
\hline
\end{tabular}
\end{center}

\begin{erwthma}
  \item Να μελετήσετε τη μονοτονία και τα ακρότατα της $f$.
  \item Ποιο είναι το ολικό μέγιστο και ποιο το ολικό ελάχιστο της $f$ στο $[-3, 5]$;
  \item Η εξίσωση $f(x) = 2$ πόσες λύσεις παρουσιάζει στο $[-3, 5]$; (Bolzano.)
  \item Να υπολογίσετε τη διαφορά $f(5) - f(-3)$ μέσω $\dint_{-3}^{5} f'(x)dx$ και το πρόσημό της.
\end{erwthma}
\end{thema}
